\documentclass[10pt,a4paper]{amsart}
\usepackage[margin=19mm]{geometry}
\usepackage{amsmath,amssymb,mathtools}
\usepackage[scr=rsfs,cal=euler]{mathalfa}
\usepackage{xfrac}
\usepackage[unicode,hidelinks,bookmarks=false]{hyperref}
\newcommand{\ie}{i.e.\ }
\newcommand{\cf}{cf.\ }
\newcommand{\resp}{respectively\ }
\newcommand{\Subsection}{\subsection}
\providecommand{\nobreakdash}{\nobreak\hbox{-}}
\setlength{\emergencystretch}{2em}
\multlinegap0pt
\usepackage[shortlabels]{enumitem}
\usepackage{amssymb}
\newtheorem{lemma}{Lemma}
\newtheorem{proposition}{Proposition}
\newcommand{\C}{C_5^{3}}
\newcommand{\Dk}{\mathsf{D}}
\newcommand{\sle}[1]{\mathsf{s}_{\le #1}}
\newcommand{\zz}{\mathsf{z}}
\title{The fourth generalized Davenport constant of $C_5^3$}
\author{Sze Chun Yiu}
\date{Source: arXiv:2609.04950v1 (CC BY 4.0)}
\begin{document}
\maketitle

\noindent\emph{Source and licence.} The fourth generalized Davenport constant of $C_5^3$. Version arXiv:2609.04950v1 (CC BY 4.0), distributed by arXiv under the Creative Commons Attribution 4.0 International licence, http://creativecommons.org/licenses/by/4.0/, as recorded in the arXiv licence metadata for this exact version and shown on the abstract page https://arxiv.org/abs/2609.04950v1. Full licence notice: \texttt{licenses/arxiv-2609.04950-CC-BY-4.0.txt}.\par\smallskip

\noindent\textit{Editorial scope.} Counted unit: the article's lemma lem:reduction (source0.tex lines 150--152). The article's complete original proof of it is delivered with it (source0.tex lines 154--156, one \texttt{proof} environment of 3 lines). No mathematics is written, repaired or re-derived: the statement and the proof are the article's own lines, in the article's own notation, and no retained line is substituted or re-flowed.  Retained uncounted as the source-local context the proof needs: lines 57--59; lines 82--87; lines 118--120; lines 134--138; lines 301--301.  Nothing was omitted from any retained span, so the package carries no omission record.  No retained line contains a citation, so the wrapper carries no bibliography and no bibliography entry.  No retained line contains a figure environment, an image-inclusion command or drawing code, so no image is delivered and the package declares no asset dependency.  Editorial formatting only, in addition: an AMS \texttt{amsart} preamble that restates the article's own declarations and macros used by the retained lines, namely the environments \texttt{lemma}, \texttt{proposition} on separate counters, together with the standard packages those lines need.  The retained environments print in this document as Lemma 1, Proposition 1 in order of appearance, whereas the article prints the same items with the numbering of its own full document.  The complete native source of the article as distributed in the arXiv e-print for this exact version is bundled unchanged as \texttt{sources/209400/source0.tex}.
\begin{equation}\label{eq:lower}
\Dk_k(\C)\ \ge\ 13+6+1+5(k-2)\ =\ 5k+10\qquad(k\ge2).
\end{equation}

\begin{lemma}\label{lem:Dk-char}
For every finite abelian group $G$ and $k\ge1$,
\[
\Dk_k(G)=\max\{\,|B| : B \text{ a zero-sum sequence over } G \text{ with } \zz(B)\le k\,\}.
\]
\end{lemma}

\begin{proposition}\label{prop:D2}
$\Dk_2(\C)=20$.
\end{proposition}

\begin{enumerate}[label=\textup{(M\arabic*)},leftmargin=3em]
\item\label{M1} $\Dk_3(\C)=25$. The lower bound is \eqref{eq:lower} at $k=3$; the archive also records the explicit witness $e_1^4\,e_2^4\,e_3^9\,(e_1+e_2)^2\,(e_1+e_3)^2\,(e_2+e_3)^3$ of length $24$ with two but not three pairwise disjoint nonempty zero-sum subsequences, checked by two separately written routines. The upper bound was obtained as follows. Suppose $S$ is a sequence of length $25$ with fewer than three pairwise disjoint nonempty zero-sum subsequences. If $A$ is any nonempty zero-sum subsequence of $S$, then $S\cdot A^{-1}$ has no two disjoint zero-sum subsequences, so $25-|A|\le\Dk_2(\C)-1=19$ by Proposition~\ref{prop:D2}, and $|A|\ge6$; by \ref{M2} below, $S$ has a zero-sum subsequence of length at most six. Hence $S=R\cdot A$ with $|A|=6$, $\sigma(A)=0$, and $R$ a sequence of length $19$ with no two disjoint zero-sum subsequences, necessarily of rank three because $\Dk(C_5^2)=9$ and $19\ge2\cdot9+1$. The archived computation enumerated all such $R$ in normalized form (support containing $e_1,e_2,e_3$ with nondecreasing multiplicities, which every rank-three $R$ can be brought to by $\mathrm{GL}(3,5)$ and a coordinate permutation), $98{,}622$ of them, then all zero-sum extensions $A$ of length six for which $R\cdot A$ has no zero-sum subsequence of length at most five (a necessary condition by the bound $|A|\ge6$ applied to $R\cdot A$), $230{,}983$ candidates in all, and found that every candidate has three pairwise disjoint nonempty zero-sum subsequences. Hence no such $S$ exists and $\Dk_3(\C)\le25$. This computation therefore rests on Proposition~\ref{prop:D2} and on \ref{M2}. Its provenance is a single pair of programs (enumeration and extension, then the three-disjoint test) with a four-input control harness; a second generator later regenerated both candidate sets with the same counts, $98{,}622$ and $230{,}983$, but the final three-disjoint test was not re-executed on them, and an earlier attempt to replay the enumeration on another machine terminated before producing a result. We have not repeated this computation.
\item\label{M2} $\sle{6}(\C)=24$; equivalently, the longest sequence over $\C$ without a nonempty zero-sum subsequence of length at most six has length $23$. The archived value comes from a single exhaustive search program over normalized sequences (rank three) together with a direct enumeration inside $C_5^2$ (rank at most two). For this paper the value was recomputed by an independently written program of a different design, which also reproduces the neighbouring values $\sle{7}(\C)=19$ and $\sle{8}(\C)=18$ recorded in the archive (Section~\ref{subsec:thresholds}).
\item\label{M3} $\sle{7}(\C)=19$, used only in the second route of Proposition~\ref{prop:D2}; same provenance and recomputation as \ref{M2}.
\end{enumerate}

\subsection{Second implementation of the threshold values}\label{subsec:thresholds}

\begin{lemma}\label{lem:reduction}
Assume \ref{M1}. If $\Dk_4(\C)=31$, then there is a zero-sum sequence over $\C$ of length $31$ with no nonempty zero-sum subsequence of length at most five.
\end{lemma}

\begin{proof}
By Lemma~\ref{lem:Dk-char} there is a zero-sum sequence $S$ of length $31$ with $\zz(S)\le4$. Let $A$ be a nonempty zero-sum subsequence of $S$ of minimum length $m$. Then $S\cdot A^{-1}$ is zero-sum, and $\zz(S\cdot A^{-1})\le3$, since four disjoint zero-sum subsequences of $S\cdot A^{-1}$ together with $A$ would give $\zz(S)\ge5$. By Lemma~\ref{lem:Dk-char} and \ref{M1}, $31-m\le\Dk_3(\C)=25$, so $m\ge6$: every nonempty zero-sum subsequence of $S$ has length at least six.
\end{proof}

\end{document}
